\section*{Chapter \ref{ch:iv1:timed-test-4} --- Timed Test 4}
\markboth{\protect\omTocMark{}{Answers: Timed Test 4}}{\protect\omTocMark{}{Answers: Timed Test 4}}

\begin{solution}{iq:iv1:timed-test-4:1}
No match has probability $\frac{8}{21}$, and exactly one pair has probability $\frac47$.

Taking the socks one at a time, the first can be anything. The second must avoid its partner: 8 of 9 remain. The third must avoid two partners: 6 of 8. The fourth must avoid three: 4 of 7. The product is $\frac89 \cdot \frac68 \cdot \frac47 = \frac{8}{21}$\index{sequential counting}.

Exactly one pair: choose the pair, in 5 ways, then 2 of the other 4 pairs, and one sock from each: $5 \times \binom42 \times 4 = 120$, out of $\binom{10}{4} = 210$. That gives $\frac47$.

Two pairs have $\binom52 = 10$ ways, or $\frac{1}{21}$, and the three cases sum to 1, a useful check.

\iqwrong{Counting ordered draws in one place and unordered ones in another, which is the usual way to get the second figure wrong.}
\end{solution}

\begin{solution}{iq:iv1:timed-test-4:2}
The probability is $\frac14$.

Given $Y + Z = s$, we have $P(X > s) = e^{-s}$. So the answer is $\mathbb{E}e^{-(Y + Z)} = (\mathbb{E}e^{-Y})^2$. And $\mathbb{E}e^{-Y} = \int_0^\infty e^{-2y}\,dy = \frac12$. The result is $\frac14$\index{memoryless property}.

A second route uses a race. Think of three exponential clocks. The first ring among $X$ and $Y$ is $Y$ with chance $\frac12$. By the memoryless property, $X$ starts anew. Then $Z$ must also beat it, again with chance $\frac12$. That gives $\frac12 \times \frac12$. A simulation confirms this.

\iqfollow{What is $P(X > 2Y)$?}{$\mathbb{E}e^{-2Y} = \int_0^\infty e^{-3y}\,dy = \frac13$.}
\end{solution}

\enlargethispage{-3\baselineskip}%
\begin{solution}{iq:iv1:timed-test-4:3}
The slope is $\frac12$, the intercept is $\frac16$, and $R^2 = \frac34$.

The means are $\bar{x} = 1$ and $\bar{y} = \frac23$. Then $S_{xx} = 1 + 0 + 1 = 2$ and $S_{xy} = (-1)(-\frac23) + 0 + (1)(\frac13) = 1$. So the slope is $\frac12$\index{least squares}. The line passes through the means, so the intercept is $\frac23 - \frac12 = \frac16$.

The total sum of squares is $\frac49 + \frac19 + \frac19 = \frac23$, and the explained part is $b^2S_{xx} = \frac12$. So $R^2 = \frac12/\frac23 = \frac34$.

The residuals are $-\frac16$, $\frac13$ and $-\frac16$. They sum to 0, and they are orthogonal to $x$, as they must be. That catches any slip.

\iqwrong{Forcing the line through the origin, which gives a slope of $\frac35$. Nothing in the question drops the intercept.}
\end{solution}

\begin{solution}{iq:iv1:timed-test-4:4}
The expected number is 6, from any square below 6.

Let $e_k$ be the expected number of rolls from square $k$. From 5, only a roll of 1 ends the game, so $e_5 = 6$.

From 4: a 1 goes to 5, a 2 finishes, and the other four rolls stay. So $e_4 = 1 + \frac16e_5 + \frac46e_4$, which gives $e_4 = 6$. By induction\index{first-step analysis}, if $e_{k + 1}, \dots, e_5$ are all 6, then from $k$ we have $6 - k$ useful rolls, of which one finishes and the rest lead to squares with value 6. Then
\[
e_k = 1 + \frac{5 - k}{6}6 + \frac{k}{6}e_k,
\]
which gives $e_k = 6$.

The reason is simple: from any square, exactly one roll finishes the game on that turn. So each roll succeeds with chance $\frac16$, whatever the square, and the wait is geometric. A simulation agrees.

\iqwrong{Computing a long recursion and never noticing the pattern. A surprising answer of 6 should prompt a search for the simple reason.}
\end{solution}

\begin{solution}{iq:iv1:timed-test-4:5}
The expected distance is $\frac{4}{\pi} \approx 1.27$.

Fix the first point. The angle $\theta$ to the second is uniform on $[0, 2\pi]$. The chord has length $2\sin(\theta/2)$\index{chord length}. So
\[
\mathbb{E} = \frac{1}{2\pi}\int_0^{2\pi}2\sin\frac{\theta}{2}\,d\theta = \frac{1}{2\pi}\left[-4\cos\frac{\theta}{2}\right]_0^{2\pi} = \frac{8}{2\pi}.
\]
A sanity check: the distance is between 0 and 2, and a chord at a right angle has length $\sqrt{2} = 1.41$. The mean is below this since the short chords, at small $\theta$ and near $2\pi$, pull it down. A simulation confirms $1.273$.

\iqwrong{Using the root mean square. The mean squared distance is exactly 2, so the root is $1.41$, which is too high. A mean of a root is not the root of a mean.}
\end{solution}

\begin{solution}{iq:iv1:timed-test-4:6}
The average is $\frac94$ rounds.

With three players, there are 27 outcomes. All the same: 3. All different: 6. So a replay has chance $\frac13$. Otherwise exactly two symbols appear. Then one winner has chance $\frac13$ (9 outcomes), and two winners also has chance $\frac13$ (9 outcomes).

With two players, a tie has chance $\frac13$, so $E_2 = \frac32$.

For three\index{first-step analysis},
\[
E_3 = 1 + \frac13E_3 + \frac13E_2 = 1 + \frac13E_3 + \frac12,
\]
which gives $\frac23E_3 = \frac32$, so $E_3 = \frac94$. A simulation of the game returns $2.25$.

\iqfollow{With four players, how likely is a replay?}{Exactly two symbols appear in $3(2^4 - 2) = 42$ of the 81 outcomes, so a replay has chance $\frac{39}{81} = \frac{13}{27}$, nearly one in two. Larger groups stall.}
\end{solution}

\begin{solution}{iq:iv1:timed-test-4:7}
The probability is $\frac{7}{10}$.

Set up states\index{first-step analysis} for the current run: $h$ after one head, $t$ after one tail, $u$ after two tails, and the start. Let each symbol stand for the chance that heads wins from that state:
\[
h = \frac12 + \frac12t, \qquad t = \frac12h + \frac12u, \qquad u = \frac12h.
\]
Then $t = \frac12h + \frac14h = \frac34h$, and $h = \frac12 + \frac38h$, which gives $h = \frac45$. So $t = \frac35$.

From the start, the first flip leads to $h$ or $t$ with equal chance: $\frac12(\frac45 + \frac35) = \frac{7}{10}$. An exact solve and a simulation agree.

So a target that is one flip shorter wins $70\%$ of the time, a little more than the raw chance of each run suggests. A limit case: against four tails in a row, the head run's chance would be higher still.

\iqwrong{Comparing $\frac14$ and $\frac18$ directly, which gives $\frac23$. This ignores that a failed attempt at one run resets differently, and that a head run resets the tail run.}
\end{solution}

\begin{solution}{iq:iv1:timed-test-4:8}
The expected number is $1 + \ln 2 \approx 1.693$.

Take logs. $E_i = -\ln U_i$ is exponential with mean 1, and the product falls below $\frac12$ when $\sum E_i$ passes $\ln 2$.

So the $E_i$ are the gaps of a Poisson process at rate 1\index{Poisson process}. The number of arrivals in $[0, \ln 2]$ is Poisson with mean $\ln 2$. The first arrival after $\ln 2$ is the stopping draw. So the number of draws is this Poisson count plus 1, whose mean is $1 + \ln 2$. A simulation gives $1.693$.

A second route: let $m(x)$ be the expected number from threshold $x$, so $m(x) = 1 + \int_x^1m(x/u)\,du$. This gives $m(x) = 1 - \ln x$, and at $\frac12$ it is the same.

\iqwrong{Using $\mathbb{E}\prod U_i = 2^{-n}$ and solving for $n = 1$. That is a statement about means, not about the first passage.}

\iqfollow{What if the target is $e^{-5}$?}{$1 + 5 = 6$ draws, by the same argument.}
\end{solution}

\begin{solution}{iq:iv1:timed-test-4:9}
Stake $\frac14$ of wealth, for a log growth of about $0.059$ per bet.

A stake of $f$ multiplies wealth by $1 + 2f$ or $1 - f$. The expected log is
\[
g(f) = \frac12\ln(1 + 2f) + \frac12\ln(1 - f).
\]
Setting the derivative to zero: $\frac{1}{1 + 2f} = \frac{1}{2(1 - f)}$, so $f = \frac14$\index{Kelly criterion}. The general form is edge over odds: $(bp - q)/b = (1 - \frac12)/2$.

Then $g = \frac12\ln 1.5 + \frac12\ln 0.75 = \frac12\ln 1.125 = 0.0589$. A direct search confirms the optimum.

Doubling to $f = \frac12$ gives $\frac12\ln(2 \times 0.5) = 0$: no growth at all. So too much leverage on a good bet leaves nothing, even with a positive expected return per bet.

\iqwrong{Staking everything, because the expected return is positive. The first loss then ends the game.}
\end{solution}

\begin{solution}{iq:iv1:timed-test-4:10}
The stronger team wins with probability about $0.710$, and equal teams play $5.81$ games on average.

Imagine all seven games are played. The team wins the series exactly when it wins at least 4 of the 7\index{extended game}. That gives
\[
\sum_{k = 4}^7\binom7k0.6^k0.4^{7 - k} = 0.2903 + 0.2613 + 0.1306 + 0.0280 = 0.710.
\]

For equal teams, the series ends after $n$ games when the winner takes game $n$ and 3 of the earlier $n - 1$. With both teams counted:
\[
P(n) = 2\binom{n - 1}{3}2^{-n}, \qquad P(4) = \frac18, \quad P(5) = \frac14, \quad P(6) = P(7) = \frac{5}{16}.
\]
The mean is $\frac{4}{8} + \frac{5}{4} + \frac{30 + 35}{16} = 5.8125$. An exact sum over games confirms both.

\iqwrong{Using $\binom{n}{4}$ for a series that ends at game $n$. The last game must be a win, so only the earlier $n - 1$ games are free. Playing all seven games avoids this, as the extra games cannot change the winner.}

\iqwhy{Extending a stopped game to a fixed length often removes the hardest part of the calculation.}
\end{solution}

\begin{solution}{iq:iv1:timed-test-4:11}
The covariance is $\frac{1225}{1296} = (\frac{35}{36})^2 \approx 0.945$.

Since $Mm$ is the product of the two dice, $\mathbb{E}Mm = 3.5^2 = 12.25$\index{covariance}.

The largest has $P(M \le k) = k^2/36$, which gives $\mathbb{E}M = \frac{161}{36}$. By the symmetry $x \to 7 - x$, $\mathbb{E}m = 7 - \frac{161}{36} = \frac{91}{36}$.

So the covariance is $12.25 - \frac{161 \times 91}{1296} = \frac{15\,876 - 14\,651}{1296} = \frac{1225}{1296}$. An enumeration over the 36 outcomes confirms it.

It is positive, as high values of one come with high values of the other. Note how the identity $Mm = XY$ avoids a table of 36 products.

\iqwrong{Assuming that $M$ and $m$ are uncorrelated because the dice are independent. They are functions of the same two dice.}
\end{solution}

\begin{solution}{iq:iv1:timed-test-4:12}
The mean is 50, the standard deviation is $\sqrt{500} \approx 22.4$, and the chance of zero is $e^{-10}$.

By the law of total expectation, $\mathbb{E}S = \mathbb{E}N\,\mathbb{E}X = 50$\index{compound Poisson}.

By the law of total variance,
\[
\operatorname{Var}S = \mathbb{E}N\operatorname{Var}X + \operatorname{Var}N(\mathbb{E}X)^2 = 10 \times 25 + 10 \times 25 = 500.
\]
For a Poisson count, this simplifies to $\lambda\mathbb{E}X^2 = 10 \times 50$. Half of the variance comes from the number of claims, and half from their sizes. A simulation agrees.

The total is 0 only with no claims, so the chance is $e^{-10} \approx 4.5 \times 10^{-5}$.

\iqwrong{Giving $\operatorname{Var}S = \lambda\operatorname{Var}X = 250$. That leaves out the uncertainty in how many claims arrive.}

\iqfollow{How does the relative risk change with $\lambda$?}{$\sqrt{\lambda\mathbb{E}X^2}/(\lambda\mathbb{E}X)$ falls as $1/\sqrt{\lambda}$. A book that is four times the size has half the relative risk.}
\end{solution}
